% Copyright 2004 by Till Tantau <tantau@users.sourceforge.net>.
%
% Modified (c) 2022 by Paul Nong-Laolam
% for ESPEC North America, INC. Web Controller V3 training distribution 
%
% In principle, this file can be redistributed and/or modified under
% the terms of the GNU Public License, version 2.
%
% However, this file is supposed to be a template to be modified
% for your own needs. For this reason, if you use this file as a
% template and not specifically distribute it as part of a another
% package/program, I grant the extra permission to freely copy and
% modify this file as you see fit and even to delete this copyright
% notice. 

\documentclass{beamer}
% Replace the \documentclass declaration above
% with the following two lines to typeset your 
% lecture notes as a handout:
%\documentclass{article}
%\usepackage{beamerarticle}


% There are many different themes available for Beamer. A comprehensive
% list with examples is given here:
% http://deic.uab.es/~iblanes/beamer_gallery/index_by_theme.html
% You can uncomment the themes below if you would like to use a different
% one:
%\usetheme{AnnArbor}
%\usetheme{Antibes}
%\usetheme{Bergen}
%cis275-ch6-w22.pdf\usetheme{Berkeley}
%\usetheme{Berlin}
%\usetheme{Boadilla}
%\usetheme{boxes}
%\usetheme{CambridgeUS}
%\usetheme{Copenhagen}
%\usetheme{Darmstadt}
%\usetheme{default}
%\usetheme{Frankfurt}
\usetheme{Goettingen}  % selected for usage
%\usetheme{Hannover}
%\usetheme{Ilmenau}
%\usetheme{JuanLesPins}
%\usetheme{Luebeck}
%\usetheme{Madrid}
%\usetheme{Malmoe}
%\usetheme{Marburg}
%\usetheme{Montpellier}
%\usetheme{PaloAlto}
%\usetheme{Pittsburgh}
%\usetheme{Rochester}
%\usetheme{Singapore}
%\usetheme{Szeged}
%\usetheme{Warsaw}

% Figures
\usepackage{graphicx}
\usepackage{colortbl}
\usepackage{verbatim} 

% explore
\usepackage{wrapfig}    % wrap text aroudn figures 
\usepackage{caption}    % set control caption for minipage env
% end explore

% menu keys
\usepackage[os=win]{menukeys}        % use key obxes
\renewmenumacro{\keys}[+]{shadowedroundedkeys}

\title{ESPEC Web Controller V3}

% A subtitle is optional and this may be deleted
\subtitle{Introducing Web Controller V3: Tech. Training}

\author{Paul Nong-Laolam}%\inst{1} } %\and S.~Another\inst{2}}
% - Give the names in the same order as the appear in the paper.
% - Use the \inst{?} command only if the authors have different
%   affiliation.

\institute[Universities of Somewhere and Elsewhere] % (optional, but mostly needed)
{
  %\inst{1}%
  ESPEC North America, INC \\ 
  4141 Central Parkway \\
  Hudsonville, Michigan 49426 \\ (USA) 
  %\and
  %\inst{2}%
  %Department of Theoretical Philosophy\\
  %University of Elsewhere
  }
% - Use the \inst command only if there are several affiliations.
% - Keep it simple, no one is interested in your street address.

\date{8/25/2022}
% - Either use conference name or its abbreviation.
% - Not really informative to the audience, more for people (including
%   yourself) who are reading the slides online

\subject{Software} % \subjectComputer Science}
% This is only inserted into the PDF information catalog. Can be left
% out. 

% If you have a file called "university-logo-filename.xxx", where xxx
% is a graphic format that can be processed by latex or pdflatex,
% resp., then you can add a logo as follows:

% \pgfdeclareimage[height=0.5cm]{university-logo}{university-logo-filename}
% \logo{\pgfuseimage{university-logo}}

% Delete this, if you do not want the table of contents to pop up at
% the beginning of each subsection:
\AtBeginSubsection[]
{
  \begin{frame}<beamer>{Outline}
    \tableofcontents[currentsection,currentsubsection]
  \end{frame}
}

% Let's get started
\begin{document}

\begin{frame}
  \titlepage
\end{frame}

\begin{frame}{ESPEC Web Controller, V3.2}
  \tableofcontents
  % You might wish to add the option [pausesections]
\end{frame}
%

%

\section{1. Introduction}

\begin{frame}{1.0: Introduction}{} 
ESPEC Web Controller V3 (EWC V3) is an embedded computer. Its hardware is based on an x86 architecture, called Up2Board (pictured below). Its operating system (OS) is powered by Debian GNU/Linux distribution.  
\begin{figure} [hbt!] 
\centering
\includegraphics[scale=.30]{figures/up2board-io-002.PNG}
\caption{EWC Up2Board with input/output ports}
\end{figure}
\end{frame}

% new slide 

\begin{frame}{1.1: ESPEC Web Controller V3}{} 
\begin{itemize}
   \item EWC V3 uses a standard Web browser to display (i.e., host) its user interface (UI) for controlling the chamber; hence, the name {\bf{Web Controller}}.  

   \item EWC V3 can operate as a standalone machine using a touchscreen or a standard display. It uses a subset of a Web browser, called Electron, to host its UI on the display.
    
   \item EWC V3 can join a network to host its UI through a Web browser, accessable through its IP address or hostname. Thanks to the ability of a Web browser to operate from any computer on the network, chamber operations can be performed via EWC V3 remotely by authorized users on the network.  

   \item No device outside the network can access EWC V3, unless the network is configured to host port forwarding to make EWC V3 visible on the Internet. 

\end{itemize}
\end{frame} 

% end of new slide 

\begin{frame}{1.2: Security Features}{} 
\begin{itemize}
% Here is another example of security issues.
   \item {\bf{Security:}} With Debian GNU/Linux continually receiving security updates, EWC V3 is very secure. 
It is also configured with read-only access to prevent intruders from exploiting the system. 

   \item {\bf{Stability:}} EWC V3 firmware update is managed by MenderIO which employs a dual-root (stepladder) technique to provide a smooth update. One root system is always in a stable state to ensure an uninterrupted operation, should the firmware update fail.

   \item {\bf{Access Privilege \& User Policy:}} EWC V3 is shipped with one user account (called admin) with complete powers to manage the system. This {\bf{admin}} account can enforce a user policy on EWC V3 and administer additional accounts with various privileges. Only authorized users can log into and control EWC V3.

\end{itemize}
\end{frame} 

\section{2. Touchscreen \& Standard Display} 

\begin{frame}{2.0: Touch-screen and Standard Display}{}
EWC V3 now supports a touchscreen or a standard display.
\begin{itemize}
   \item {\bf{Touchscreen:}} 
    With a touchscreen (or a standard display) attached to its display port, EWC V3 automatically displays its UI for chamber control and operation. 
    \begin{figure} [hbt!] 
      \centering
      \includegraphics[scale=0.30]{figures/touchscreen-panel-001a.PNG}  
      \caption{EWC V3 with its Touchscreen (HMI)}
    \end{figure}
\end{itemize}  
\end{frame}            
%

\begin{frame}{2.1: Touchscreen and Standard Display}{}
Availability of display options and advantages:
\begin{itemize}
   \item {\bf{Touchscreen:}} 
         \begin{enumerate}
             \item Typhoon chambers are shipped with the MIMO touchscreen to provide the UI of EWC V3 as its dedicated HMI. 
             \item A custom chamber (with F4T) will be equipped with the MIMO 
                                     touchscreen to provide the UI of EWC V3 as its dedicated HMI.
         \end{enumerate} 
   \item {\bf{Standard Monitor:}} 
          \begin{itemize}
               \item With a standard display, mouse and keyboard attached, EWC V3 can be operated 
                     as a standalone device to control the chamber. 
          \end{itemize} 
   \item {\bf{Advantage of Touchscreen or Standard Display:}} 
          \begin{itemize}
               \item If EWC V3 cannot be part of a network, it can operate as a standalone device.  
               \item Troubleshooting a network connection is easy. 
          \end{itemize} 
\end{itemize}  
\end{frame}    
%

\section{3. Cloud Updating} 
\begin{frame}{3.0: Cloud Updating}{}
EWC V3 uses chamber S/N as its registration number to receive firmware update through a cloud service provider, called MenderIO.
    \begin{figure} [hbt!] 
      \centering
%      \includegraphics[scale=.30]{figures/menderio-cloud-001.PNG}  
      \includegraphics[scale=0.320]{figures/mender-cloud-service-002.PNG}
      \caption{EWC subscribes to MenderIO cloud service}
    \end{figure}
\end{frame} 

\begin{frame}{3.1: Firmware Update: Online or Offline}{} 
Customer has complete control over firmware update method: Offline or Online through a cloud service. 
\begin{enumerate}
\item Using a stepladder technique, update package is pushed onto the next available root partition (see diagram). 

    \begin{figure} [hbt!] 
      \centering
%      \includegraphics[scale=.30]{figures/menderio-cloud-001.PNG}  
      \includegraphics[scale=.60]{figures/stepladder-001.PNG}
      \caption{Stepladder technique of EWC V3 update}
    \end{figure}

\item If the update process is successful, the system will boot and run on this new firmware.   
\item If the update process failed, the system rolls back to the previous EWC V3 known to be in the stable state. 
\end{enumerate} 
\end{frame} 

\begin{frame}{3.2: Firmware Update Subscription}{} 
Firmware update is provided through a yearly update subscription. Customers purchase the subscription 
to receive the updates. 
\begin{itemize}
   \item {\bf{Online Update}}: \\
       For EWC V3 that connects to the Internet, MenderIO cloud updating service handles the update process. 
       Automatic update will take place when the chamber (connected EWC V3) is in a Standby mode.   
   \item {\bf{Offline Update}}:   \\
       For EWC V3 that does not connect to the Internet, the owner is provided with a link to access their account
       to download an update package to perform the update procedure on their system. Refer to User's
       Manual for detail on Firmware Update procedure. 
\end{itemize} 
\end{frame} 

\section{4. User Interface: Display Options}
\begin{frame}{4.0: Display Options}{}

\begin{minipage}[c]{0.50\textwidth}
\vspace{0.125in} 

\begin{itemize} 
   \item On a touchscreen (or a standard display), EWC V3 provides an on-screen keyboard and
options for attaching a USB drive, taking screenshots, downloading data, setting language or date/time, etc.  
\end{itemize} 
\end{minipage}
\hfill
\begin{minipage}[c]{0.450\textwidth}

%\begin{figure} [hbt!] 
%      \centering
%      \includegraphics[scale=.15]{figures/onscreen-login-001a.PNG} 
      \includegraphics[width=\textwidth]{figures/onscreen-login-001a.PNG}  
%      \caption{EWC V3 log-in screen on a touchscreen}
      \captionof{figure}{EWC V3 log-in screen on a touchscreen}
%\end{figure}
\end{minipage} 
\vspace{0.2in} 

\begin{itemize} 
    \item EWC V3 UI on a Web browser (PC). 
\end{itemize} 

\begin{figure} [hbt!] 
      \centering
      \includegraphics[scale=.15]{figures/desktop-login-001.PNG}  
      \caption{EWC V3 log-in screen on a Web browser}
\end{figure}
\end{frame} 

\begin{frame}{4.1: User Interface}{}
EWC V3 got a new makeover; a completely new UI.  
    \begin{figure} [hbt!] 
      \centering
      \includegraphics[scale=.18]{figures/homepage-001.PNG}  
      \caption{EWC V3 Overview page}
    \end{figure}

\begin{itemize}  
   \item {\bf{Menu bar:}} Access to all the system's menus to control and operate the chamber 
   \item {\bf{Status bar:}} Provides a quick summary of the current status and condition of the chamber. 
   \item {\bf{User Account:}} One user account, called admin. 
   \item {\bf{User's Manual:}} Has embedded User's Manual   
\end{itemize} 
\end{frame}

\section{5. Chamber Support \& Interface, Network Configuration} 
\begin{frame}{5.0: Communication Interface \& Configuration}{}

ESPEC Web Controller V3 currently supports these PLCs: 
\begin{itemize}
   \item {\bf{T-Series Chamber:}} Allen Bradley CompatLogix, ControlLogix and Micro8xx series.\\ 
                                  {\it{Note}}: A touchscreen is the default HMI for EWC V3 shipped with T-series chambers. 
   \item {\bf{ESPEC Chambers:}} ESPEC P300, SCP220, ES102 
   \item {\bf{Watlow:}} F4T and F4 \\
                        {\it{Note}}: A touchscreen is the default HMI for EWC V3 shipped with a custom chamber with F4T.  
\end{itemize}

Default Communication between EWC V3 and Chamber:
\begin{enumerate}
   \item {\bf{TCP/IP Interface:}} T-series and F4T 
   \item {\bf{Serial Interface:}} 
          \begin{itemize}
               \item {\bf{RS-485}}: P300
               \item {\bf{RS-232}}: SCP220, ES102, F4
          \end{itemize}  
\end{enumerate}
\end{frame}

\begin{frame}{5.1: DHCP \& Static Nework Configuration}{}
\begin{itemize} 
   \item DHCP: IP address handled by DHCP server
       \begin{figure} [hbt!] 
         \centering
         \includegraphics[scale=.14]{figures/1672855610-dhcp-setup-figure.PNG}  
         \caption{DHCP network diagram}
       \end{figure}

   \item Static Network: IP address configured manually
      \begin{figure} [hbt!] 
        \centering
        \includegraphics[scale=.18]{figures/2789777267-typhoon-static-setup-v2.PNG}  
        \caption{Static network diagram}
      \end{figure}
\end{itemize} 
\end{frame}

\begin{frame}{5.2: TCP Configuration}{}
\begin{minipage}[c]{0.50\textwidth}
{\bf{Internal Network:}} 

Typhoon and F4T chambers use an internal network protocol with unique IP address to
communicate with EWC V3. This internal network is configured by EWC V3 by default, as illustrated in the figure (right).    
\end{minipage} 
\hfill
\begin{minipage}[c]{0.450\textwidth}
%      \begin{figure} [hbt!] 
%        \centering
%        \includegraphics[scale=.20]{figures/typhoon-network-cfg-001.PNG}  
%        \includegraphics[width=\textwidth]{figures/typhoon-network-cfg-002.PNG}  
%        \captionof{figure}{Typhoon network config}
%        \caption{Example of Typhoon internal network setting}
%      \end{figure}

%      \begin{figure} [hbt!] 
%        \centering
%        \includegraphics[scale=.20]{figures/f4t-network-cfg-001.PNG}  
        \includegraphics[width=\textwidth]{figures/f4t-network-cfg-002.PNG} 
%        \caption{Example of F4T internal network setting}
        \captionof{figure}{F4T network config}
%      \end{figure}
\end{minipage}

\begin{minipage}[c]{0.50\textwidth}
{\bf{External Network:}} 

EWC V3 uses a predefined Ethernet port on its hardware to connect to the customer main network as depicted in the diagram on the right.  
\end{minipage}
\hfill
\begin{minipage}[c]{0.450\textwidth}
%      \begin{figure} [hbt!] 
%        \centering
%        \includegraphics[scale=.20]{figures/f4t-network-cfg-001.PNG}  
        \includegraphics[width=\textwidth]{figures/up2board-io-002.PNG} 
%        \caption{Example of F4T internal network setting}
        \captionof{figure}{Ethernet ports on Up2board}
%      \end{figure}
\end{minipage}
\end{frame}

\begin{frame}{5.3: TCP: Chamber Communication Interface}{}

For F4T chambers, there is an additional option on chamber communication interface, through a WAN network, where the F4T and EWC V3 both connect to the customer network (instead of to each other as described above). In this setup, the customer network will have to supply two IP addresses: one for EWC V3 and one for F4T. Thus, the default configuration is more practical, since it still allows customer to access the F4T with Watlow Composer via the IP address of EWC V3. 

%      \begin{figure} [hbt!] 
%        \centering
%        \includegraphics[scale=.20]{figures/typhoon-network-cfg-001.PNG}  
%        \includegraphics[width=\textwidth]{figures/typhoon-network-cfg-002.PNG}  
%        \captionof{figure}{Typhoon network config}
%        \caption{Example of Typhoon internal network setting}
%      \end{figure}

      \begin{figure} [hbt!] 
        \centering
        \includegraphics[scale=.50]{figures/f4t-wan-cfg-001.PNG}  
%        \includegraphics[width=\textwidth]{figures/f4t-wan-cfg-001.PNG} 
%        \caption{Example of F4T internal network setting}
        \captionof{figure}{F4T WAN config via customer network}
      \end{figure}
\end{frame}

\begin{frame}{5.4: Connecting to F4T}{} 
EWC V3 and F4T communicate through TCP/IP interface. Confirm that
connections are set up are follows: 

    \begin{figure} [hbt!] 
      \centering
      \includegraphics[scale=.30]{figures/up2board-to-f4t.PNG}  \\
      \caption{F4T and Web Controller connection}
    \end{figure}

Step-by-step procedure for configuring F4T internal network will be provided.  
\end{frame}

\begin{frame}{5.5: Serial Interface Configuration}{}

Serial interface communication is dictated by the type of controller.

  \begin{itemize}
     \item {\bf{RS-485}}: The default serial interface for P300 with baud rate of 19200. However, RS-232 can also be used. 
     \item {\bf{RS-232}}: The default serial interface for SCP220, ES102 and F4 with baud rate of 9600. 
  \end{itemize}  

\begin{figure} [hbt!] 
\centering
\includegraphics[scale=.25]{figures/up2board-io-002.PNG}
\caption{EWC V3 Up2board with Serial Interface}
\end{figure}

\end{frame}

\begin{frame}{5.6: P300 with RS-485}{} 
RS-485 is the default communication interface for P300. Correct pin configuration on the Up2Board required for RS-485 (or RS-232): 

    \begin{figure} [hbt!] 
      \centering
      \includegraphics[scale=.35]{figures/p300-jumper-setting-0a.PNG}  \\
      \caption{RS-485 jumper setting on Up2Board for P300 communication}
    \end{figure}

\end{frame}

\begin{frame}{5.7: Connecting to SCP220/P300/F4}{} 
RS-232 is the standard communication interface for SCP220 and F4. P300 can its 3-pin RS-232C for communication. Correct pin configuration on the Up2Board is required for RS-232: 

    \begin{figure} [hbt!] 
      \centering
      \includegraphics[scale=.37]{figures/scp220-jumper-setting-0a.PNG}  \\
      \caption{RS-232 jumper setting on Up2Board for SCP220/F4 and P300 (3-pin) COMM}
    \end{figure}

\end{frame}
\section{6. U-Web Kit} 
\begin{frame}{6.0: U-Web Kit Support}{} 
EWC V3 can be sold as a Universal Web Kit to connect and control the existing chamber. 

    \begin{figure} [hbt!] 
      \centering
      \includegraphics[scale=.260]{figures/u-web-kit-up2board-001.PNG}  
      \caption{U-Web Kit with USB-to-Serial interface cable}
    \end{figure}
{\bf{Advantages:}} 
\begin{itemize}
\item U-Web kit is an excellent option for adding EWC V3 to the existing chamber for control and operation. 
\item Setup is generally made simple with its quick guide and preconfiguration. 
\end{itemize} 

\end{frame}

\section{7. EWC V3 Features \& Operation}

\begin{frame}{7.0: The Menu Bar of EWC V3}{}
\begin{minipage}[c]{0.50\textwidth}
All features of EWC V3 are accessable and controllable through its nine menus in the menu bar (remains fixed on the left). 

\vspace{0.125in} 
The {\bf{Constant}}, {\bf{Program}}, {\bf{Start/Stop}} and {\bf{Settings}} menus are generally linked to the user's privileges controlled by the EWC V3 administrator. 

\vspace{0.125in}
These menus are grayed out to those users without read/write privileges. 
\end{minipage}
\hfill
\begin{minipage}[c]{0.45\textwidth}
%    \begin{figure} [hbt!] 
%      \centering
%      \includegraphics[scale=.50]{figures/main-menu-001.PNG}  
      \includegraphics[width=\textwidth]{figures/main-menu-002.PNG} 
%      \caption{Main menu of EWC V3}
      \captionof{figure}{EWC V3 main menu}
%    \end{figure}

%      \includegraphics[width=\textwidth]{figures/main-menu-001.PNG} 
%      \caption{Main menu of EWC V3}
%      \captionof{figure}{EWC V3 main menu}
\end{minipage}
\end{frame}

\begin{frame}{7.1: EWC V3 Main Menu \& Features}{} 

\begin{enumerate}
   \item {\bf{User}}: login/logout menu.
   \item {\bf{Overview}}: Quick summary of the chamber operating status and conditions.
   \item {\bf{Trend}}: A scatter plot of the data points collected by the chamber.
   \item {\bf{History}}: Displays logs of the chamber operation and its statistics.
   \item {\bf{Constant}}: Setting or re-setting new constant values.
   \item {\bf{Program}}: Construct program instructions, preview program output; program file manipulation, and more.
   \item {\bf{Start Stop}}: Control the chamber operating mode.
   \item {\bf{Setting}}: Settings of various features and operations.
   \item {\bf{About}}: Information about terms and agreement of EWC V3, online support and embedded operation manual.
\end{enumerate} 
\end{frame}

\section{8. Hands-On: Demo} 

\begin{frame}{8.0: Using EWC V3, Hands-On}{}
\begin{itemize}
   \item Quick demo of EWC V3 operation and features

    \begin{figure} [hbt!] 
      \centering
      \includegraphics[scale=.150]{figures/prog-002.PNG}  
      \caption{Program/Preview menu of EWC V3}
    \end{figure}

    \begin{figure} [hbt!] 
      \centering
      \includegraphics[scale=.150]{figures/server-001.PNG}  
      \caption{Server menu of EWC V3}
    \end{figure}

   \item Questions \& Answers
\end{itemize}
\end{frame}


% All of the following is optional and typically not needed. 
\appendix
%\section<presentation>*{\appendixname}
%\subsection<presentation>*{References}
\section<presentation>*{References}

\begin{frame}[allowframebreaks]
  \frametitle<presentation>{References}
    
  \begin{thebibliography}{10}
    
  \beamertemplatebookbibitems
  % Start with overview books.

  \bibitem{ESPECWebManual2022}
    ESPEC.~Web Controller V3.
    \newblock {\em User Manual: ESPEC Web Controller, V3}, wiki/bitbucket.com, 2022. \\
    {\url{https://bitbucket.org/especnorthamerica/especweb/wiki/Web_Controller_V3_Home}}

  \bibitem{ESPECWebManual2022}
    ESPEC.~Web Controller V3.
    \newblock {\em User Manual} [ESPEC Web Controller, V3]. Retrievable and viewable
    on the menu bar under {\sl{About}}, 2022. \\
    {\url{http://especSERIAL\#.local/\#/about}}

  \bibitem{ESPECWebConfig2022}
    ESPEC.~Production Configuration.
    \newblock {\em ESPEC Web Controller V3, Factory Config}, ESPEC Web Controller V3, Internal Departmental Resource, 2022. \\
    {\url{https://bitbucket.org/especnorthamerica/espec-web-vue/wiki/Factory_Setup_Procedure}}
    %\newblock MCC, 2020.
  

%  \beamertemplatearticlebibitems
%  \bibitem{UbuntuUnity2017}
%     Unity desktop ended.
%     \newblock Failed to popualrize or market Unity as a solution to all Ubuntu products, Unity project will be ended in April 2017, said Canonical boss Mark Shuttleworth.  Retrieved from {\url{https://www.bbc.com/news/technology-39490848}}. 


  %  S.~Someone.
  %  \newblock On this and that.
  %  \newblock {\em Journal of This and That}, 2(1):50--100,
  %  2000.
  \end{thebibliography}
\end{frame}

\end{document}
